% =====================================================================
%  说明书
%
%  \p{...}             一个自然段，默认按官方空白表格不显示段号；需要
%                      校阅稿段号时，在 cnipa 选项中启用 numbered-paragraphs。
%  \cnsection{...}     一级小节：技术领域 / 背景技术 / 发明内容 / …
%  \cnsubsection{...}  发明内容下的一、二、三
%  \cnexample{...}     实施例标题
%
%  附图标记 100–111 必须与 figures/ 下的五幅图完全一致：说明书未提及的
%  标记不得出现在图中，图中出现的标记也不得在说明书里缺席。
% =====================================================================
\documentclass[12pt,UTF8,fontset=none]{ctexart}
\usepackage[description]{cnipa}

\begin{document}
% 官方表格的“说明书”在页眉；正文第一行是名称，二者互不替代。
\inventiontitle{一种大语言模型智能体的工具调用控制方法及系统}

\cnsection{技术领域}

\p{本发明涉及计算机软件技术领域，具体涉及一种由大语言模型驱动的智能体在其执行过程中对工具调用进行控制的方法及系统，尤其涉及对不产生新观察的重复工具调用的识别与阻断。}

\cnsection{背景技术}

\p{在由大语言模型驱动的智能体运行时中，一个用户回合通常由多次“模型生成工具调用—运行时执行该工具调用—将工具结果写回会话上下文—模型依据会话上下文生成下一次工具调用”的迭代构成。每一次迭代都要将全部会话上下文重新提交给所述大语言模型，因而都会消耗令牌、算力以及对外部服务的网络请求配额。}

\p{在上述迭代中，工具调用的次数并不等于任务的推进程度。若干次措辞不同的工具调用可能反复取回同一个页面、同一条运行状态、同一个空结果或同一种失败。此时会话上下文持续增长、消耗持续发生，而可用信息不再增加。若无相应的控制手段，该过程可能持续到回合预算耗尽为止。}

\p{现有的控制手段主要存在以下不足：}

\p{其一，以工具调用的重复次数作为唯一判据。此类手段对同一调用连续出现预设次数后加以抑制。但在智能体场景中，对运行状态的轮询本身就要求以相同入参重复调用，此类手段会将正当的轮询误判为无效重复；同时，只要模型对入参作出细微改写（例如改变检索词的措辞、改变分页参数而落回同一页内容），重复即被绕过，而取回的观察并未改变。}

\p{其二，工具输出中普遍含有时间戳、耗时、请求标识、链路追踪标识与心跳等随每次调用而变化的字段。若直接对工具输出整体判重，内容未发生变化的运行状态会因时间推移而每次都呈现为不同的结果，导致判重失效。}

\p{其三，为缓解上下文预算压力，运行时会将体积较大的工具输出移出会话上下文、存入本地文件，并在会话上下文中仅保留指向该文件的引用。由于该文件的存储路径中通常含有本次工具调用的调用标识以保证并发写入时文件名互不冲突，同一份内容在不同次调用下必然落在不同的路径上。此时若以会话上下文中可见的内容（即该路径）作为判重依据，同一份内容会被反复认定为新结果，外置本身即使判重失效，二者不能同时成立。}

\p{其四，当会话上下文因超出预算而被压缩时，用于判重的中间状态随原始消息一并丢失。压缩之后的会话无从得知其在压缩之前已经处于反复取回同一结果的状态，因而会重新开始同样的重复过程。}

\p{因此，需要一种能够依据工具输出的实际内容判定“本次调用是否带来新观察”，能够在大体积输出被移出会话上下文后继续成立，并且能够跨越上下文压缩而不丢失的工具调用控制手段。}

\cnsection{发明内容}

\cnsubsection{一、本发明要解决的技术问题}

\p{本发明要解决的技术问题是：在大语言模型智能体的一个用户回合内，反复发起不产生新观察的工具调用持续消耗令牌、算力与外部网络请求配额；以调用是否重复为唯一判据的抑制手段对轮询类工具误判、对入参微变的同义调用漏判；工具输出中的易变字段使内容未变的结果被判定为新结果；大体积工具输出被外置存储后其存储路径随调用标识而变，使判重随外置一并失效；以及上下文压缩导致判重所依赖的状态量丢失。}

\cnsubsection{二、本发明采用的技术方案}

\p{为解决上述技术问题，本发明采用的技术方案是，提供一种大语言模型智能体的工具调用控制方法，所述方法由智能体运行时执行，所述智能体运行时在一个用户回合内迭代地接收大语言模型生成的工具调用、执行所述工具调用以得到工具结果、将所述工具结果写入会话上下文并由所述大语言模型依据所述会话上下文生成下一次工具调用；所述方法还包括：}

\p{对本用户回合内每一次已完成的工具调用，分别生成调用指纹与证据指纹，并写入与该次工具调用对应的观察记录；所述调用指纹由工具名称与经规范化处理的调用入参生成，所述证据指纹由工具输出的规范化表示生成。所述观察记录至少保存调用指纹、证据指纹或无证据标识、规范化规则版本、调用标识和工具结果完成顺序，从而使账本状态具有可重放的输入。}

\p{生成所述证据指纹时，先对所述工具输出进行确定性的结构规范化处理并剔除预设的易变字段，再由结构规范化处理后的结果生成所述证据指纹；所述结构规范化处理包括将所述工具输出解析为结构化对象并按键名排序，所述易变字段包括时间戳字段、耗时字段、请求标识字段、链路追踪标识字段与心跳字段中的至少一项，使得内容未发生变化的工具输出生成相同的证据指纹。规范化规则以版本标识写入所述观察记录，规则升级时能够区分不同版本生成的证据身份。}

\p{对工具输出的长度超过外置阈值的工具调用，将该工具输出存储为产物文件，并在所述会话上下文中以引用式结果替换该工具输出，所述引用式结果包含所述产物文件的存储路径、有限预览与该工具输出的内容摘要值，所述存储路径中包含该次工具调用的调用标识；对于所述引用式结果，以其中的所述内容摘要值而非所述存储路径、有限预览或调用标识作为证据身份依据。由此，原始工具输出、引用式结果和后续恢复表示只要指向同一内容，其证据指纹保持一致。}

\p{维护本用户回合的调用指纹计数、证据指纹集合、同结果计数与无新证据连续计数；当证据指纹未出现在所述证据指纹集合中时，将其加入所述证据指纹集合并将所述无新证据连续计数清零；当证据指纹已出现在所述证据指纹集合中、或所述工具输出经所述结构规范化处理后为空时，将所述无新证据连续计数累加。所述账本将工具执行产生的外部观察与运行时合成的控制结果区分保存，合成控制结果不改变证据新颖性状态。}

\p{当所述无新证据连续计数达到第一阈值时，向所述工具结果中注入策略变更指示，所述策略变更指示指明下一次动作应当改变假设、来源、目标或机制。所述策略变更指示是运行时生成的内部控制结果，用于改变后续工具调用的候选空间，不作为外部证据。}

\p{在执行一次新的工具调用之前，先生成该次工具调用的调用指纹，并在所述无新证据连续计数达到第二门控条件、该调用指纹对应的调用指纹计数不小于第三门控条件、且该调用指纹对应的同结果计数不小于第四门控条件三者同时满足时，在执行前拒绝该次工具调用；所述第二门控条件的强度高于触发策略变更指示的第一门控条件。对属于预设轮询类工具集合的工具调用，所述第二门控条件根据工具能力描述取为非轮询类工具对应条件的整数倍。}

\p{当所述会话上下文被压缩、恢复或重新加载时，从持久化的观察记录，或者从所述会话上下文中已存储的工具调用与工具结果的配对出发，按照观察记录中记载的规范化规则版本和相同的确定性规则重放并重建所述调用指纹计数、所述证据指纹集合、所述同结果计数与所述无新证据连续计数。所述重放不重新调用所述大语言模型或外部工具，并保持重放前后的证据新颖性判定关系一致。}

\p{本发明还提供一种实施上述方法的系统，以及一种存储有实现上述方法的计算机程序的计算机可读存储介质。}

\cnsubsection{三、本发明的有益效果}

\p{本发明的有益效果是：}

\p{由于判定依据取自工具输出的实际内容而非工具调用的次数，措辞不同但取回同一观察的多次工具调用被识别为同一份证据，从而在空转发生时即可被察觉，而不必等到回合预算耗尽。}

\p{由于在生成证据指纹之前对工具输出作结构规范化处理并剔除易变字段，内容未发生变化的运行状态不会因时间推移、请求标识变化或心跳字段更新而被误判为新结果，判重在含有此类字段的工具输出上同样成立。}

\p{由于证据身份与输出表示解耦，对被外置存储的工具输出以其内容摘要值而非含有调用标识的存储路径、有限预览或调用标识作为证据身份，同一份内容无论被写入何种文件名、截取为何种预览或在压缩后以何种引用表示出现，其证据指纹保持不变；因而上下文预算优化与工具调用循环检测可以同时启用，前者不再以损害后者为代价。}

\p{由于阻断需要无新证据连续计数、调用指纹计数与同结果计数三者同时达到各自阈值，且对轮询类工具将其中的连续计数阈值成倍放大，正当的重试与状态轮询不被误拒；而已被证明既重复发起又重复返回同一结果的工具调用在执行之前即被拒绝，其对应的算力、令牌与外部网络请求因此不再发生。}

\p{由于所述观察记录保存了规范化规则版本和工具结果类型，所述各计数与集合可从观察记录或会话上下文中已存储的工具调用与工具结果配对按相同规则确定性重放，上下文压缩、会话恢复或进程重启不导致判重能力丢失，长回合在恢复之后仍能察觉其处于重复状态。}

\p{由于所述各计数与集合均由确定性的规范化与摘要运算得到，不依赖对模型输出的再次推理，判定过程可复现且不引入额外的模型调用开销。}

\cnsection{附图说明}

\p{图 1 是本发明所述工具调用控制系统的结构示意图；}

\p{图 2 是本发明实施例中对一次已完成的工具调用生成双重指纹并更新账本的流程图；}

\p{图 3 是本发明实施例中在执行一次新的工具调用之前进行阻断判定的流程图；}

\p{图 4 是本发明实施例中大体积工具输出的外置存储与证据身份桥接的示意图；}

\p{图 5 是本发明实施例中会话上下文被压缩时重建账本的流程图。}

\p{图中：100—智能体运行时；101—模型接口单元；102—工具执行单元；103—指纹生成单元；104—证据账本单元；105—循环控制单元；106—产物存储单元；107—会话上下文存储单元；108—上下文压缩单元；109—事件记录单元；110—大语言模型；111—工具集合。}

\cnsection{具体实施方式}

\p{下面结合附图和实施例对本发明作进一步详细描述。以下实施例用于说明本发明，但不用来限制本发明的范围。}

\cnexample{实施例一：系统结构与总体流程}

\p{参照图 1，本实施例提供的工具调用控制系统包括智能体运行时（100）。所述智能体运行时 100 与位于其外部的大语言模型（110）以及工具集合（111）通信。所述智能体运行时 100 包括模型接口单元（101）、工具执行单元（102）、指纹生成单元（103）、证据账本单元（104）、循环控制单元（105）、产物存储单元（106）、会话上下文存储单元（107）、上下文压缩单元（108）与事件记录单元（109）。}

\p{所述会话上下文存储单元 107 保存本用户回合的会话上下文。所述会话上下文由消息序列构成，每条消息具有角色与内容块序列，所述内容块包括文本块、工具调用块与工具结果块；所述工具调用块含有调用标识、工具名称与调用入参，所述工具结果块含有与之对应的调用标识、工具名称、工具输出与失败标记。}

\p{所述模型接口单元 101 将所述会话上下文提交给所述大语言模型 110，并接收所述大语言模型 110 生成的工具调用。所述工具执行单元 102 在所述工具集合 111 中执行所述工具调用并得到工具输出。所述指纹生成单元 103 依据所述工具名称、所述调用入参与所述工具输出生成调用指纹与证据指纹。所述证据账本单元 104 维护本用户回合的调用指纹计数、证据指纹集合、同结果计数与无新证据连续计数，并将每一次已完成调用的结果写成观察记录。所述观察记录保存规范化规则版本、调用标识、完成顺序、表示类型以及工具结果是否为运行时合成结果的标记。所述循环控制单元 105 依据所述证据账本单元 104 的状态，在所述工具结果中注入策略变更指示或在执行之前拒绝工具调用。所述产物存储单元 106 将长度超过外置阈值的工具输出存储为产物文件，并在引用式结果中保留内容摘要值。所述上下文压缩单元 108 在所述会话上下文超出预算时对其进行压缩，并触发所述证据账本单元 104 按观察记录重放重建。所述事件记录单元 109 输出不含工具输出本身的元数据事件。}

\p{在本实施例中，证据指纹不是会话中某一种文字表示的地址，而是工具观察的表示无关身份。原始输出、包含调用标识的产物路径、有限预览以及恢复时的引用式结果可以不同，但只要引用式结果携带同一内容摘要值，所述证据账本单元 104 即将其映射为同一证据指纹。对于不携带内容摘要值的恢复表示，所述证据账本单元 104 从持久化观察记录取回该身份，而不根据路径或预览重新生成身份。}

\p{在一个实施方式中，所述表示变换单元由所述产物存储单元 106 实现，用于将原始工具输出转换为产物文件和引用式结果；所述观察记录存储单元由所述证据账本单元 104 与所述会话上下文存储单元 107 共同实现，用于保存所述观察记录；所述上下文恢复单元由所述上下文压缩单元 108 实现，用于在压缩、恢复或重新加载时触发观察记录的确定性重放。上述单元可以由同一处理器上的程序模块实现，也可以由不同的处理器或存储介质实现。}

\p{在一个实施方式中，所述证据账本单元 104 的状态在每一个新的用户回合开始时被重置，从而本用户回合的判定不受此前回合的影响。}

\p{在一个实施方式中，所述智能体运行时 100 具有第一模式、第二模式与第三模式，可通过运行环境变量配置。所述第一模式下不执行指纹生成、账本更新、指示注入与执行前拒绝；所述第二模式下执行指纹生成、账本更新、指示注入与账本重建，但不执行执行前拒绝；所述第三模式下上述各项全部执行。所述第二模式可用于在不改变既有执行行为的前提下观察判定结果。在一个优选的实施方式中，所述第三模式为默认模式。}

\cnexample{实施例二：双重指纹的生成}

\p{参照图 2，对一次已完成的工具调用，所述智能体运行时 100 执行如下步骤。}

\p{步骤 S201，所述指纹生成单元 103 生成该次工具调用的调用指纹。具体地，先尝试将所述调用入参解析为结构化对象；若解析成功，则对该结构化对象作结构规范化处理，即递归地按键名升序重排对象的键值对、保持数组元素的原有次序，并按统一格式序列化为字符串；若解析失败，则将所述调用入参中的连续空白字符合并为单个空格。随后，将转换为小写的所述工具名称与上述结果按分隔符拼接，对拼接结果计算摘要值，所得摘要值即为所述调用指纹。所述结构规范化处理规则具有规则版本标识，该规则版本与调用指纹一并写入观察记录。}

\p{步骤 S202，所述指纹生成单元 103 对所述工具输出作预处理，剥离运行时在此前迭代中追加到所述工具输出尾部的内部提示文本，并去除首尾空白。}

\p{步骤 S203，判断预处理后的所述工具输出是否为空。若为空，则该次工具调用不产生证据指纹，判定结果为“空”，转入步骤 S207。}

\p{步骤 S204，判断预处理后的所述工具输出是否为可解析的结构化对象且符合引用式结果的形式。所述引用式结果的形式为：含有取值为预设标记的状态字段、含有非空的存储路径字段、且含有长度与取值符合摘要值格式的内容摘要值字段。若符合，则以预设前缀与该内容摘要值字段的取值拼接所得的字符串作为规范化结果，转入步骤 S206。该步骤的作用在实施例四中详述。}

\p{步骤 S205，若不符合引用式结果的形式，则区分下述两种情形生成规范化结果：}

\p{情形一，该次工具调用被标记为失败。此时从所述工具输出中提取失败描述文本；在所述工具输出为结构化对象时，依次在预设的失败字段名中查找失败描述文本，并对结构化的单元格执行结果优先提取其中的异常名称与异常取值，从而避免所有失败因其序列化外层相同而收敛到同一个描述。随后，将所述失败描述文本中的每一段连续数字串替换为单个统一占位符、将全部字符转换为小写、将连续空白字符合并为单个空格，并截断至预设长度上限。由此，同一类失败不会因其中的行号、偏移量、进程号或重试序号不同而被判定为不同的证据。在一个实施方式中，若无法提取到失败描述文本，则退回到对所述工具输出整体合并空白字符的处理。}

\p{情形二，该次工具调用未被标记为失败。此时若所述工具输出可解析为结构化对象，则对其作所述结构规范化处理；在该处理过程中，键名属于预设易变字段集合的键值对被整体剔除，且当某一键值对的键名属于所述预设易变字段集合时，其取值不参与序列化。在一个实施方式中，所述预设易变字段集合包括表示耗时的字段名、表示持续时长的字段名、表示时间戳的字段名、表示更新时刻的字段名、表示开始时刻的字段名、表示结束时刻的字段名、表示请求标识的字段名、表示链路追踪标识的字段名与表示心跳的字段名，并同时覆盖上述字段名的不同书写形式。若所述工具输出不可解析为结构化对象，则将其连续空白字符合并为单个空格。}

\p{步骤 S206，判断所述规范化结果是否为空。若剔除易变字段后所述规范化结果为空，则该次工具调用不产生证据指纹，判定结果为“空”；否则，将转换为小写的所述工具名称与所述规范化结果按分隔符拼接并计算摘要值，所得摘要值即为所述证据指纹。}

\p{步骤 S207，所述证据账本单元 104 更新账本，并写入该次工具调用的观察记录。具体地，将该次工具调用的调用指纹对应的调用指纹计数加一，将已完成工具调用的总次数加一，并在观察记录中写入调用标识、工具结果完成顺序、规范化规则版本、工具输出的原始内容摘要值、规范化内容摘要值、证据指纹或无证据标识、表示类型以及工具结果是否为运行时合成结果的标记；随后按下述规则之一处理：}

\p{规则一，若该次工具调用不产生证据指纹，则将空结果次数加一，并将所述无新证据连续计数加一，判定结果为“空”；}

\p{规则二，若所述证据指纹未出现在所述证据指纹集合中，则将其加入所述证据指纹集合，将所述无新证据连续计数清零，并清除此前记录的指示注入位置，判定结果为“新”；}

\p{规则三，若所述证据指纹已出现在所述证据指纹集合中，则将重复次数加一，并将所述无新证据连续计数加一，判定结果为“重复”。}

\p{步骤 S208，所述证据账本单元 104 更新同结果计数。具体地，以该次工具调用的调用指纹为键，取出此前记录的结果标识与计数；若该次工具调用的证据指纹与所记录的结果标识相同，则将该计数加一；否则以该次工具调用的证据指纹替换所记录的结果标识并将该计数置为一。在一个实施方式中，不产生证据指纹的情形以预设的固定标识参与上述比较。}

\p{步骤 S209，所述事件记录单元 109 输出一条元数据事件，所述元数据事件包含所述工具名称、所述判定结果、所述证据指纹的前若干位截断值、所述无新证据连续计数、所述证据指纹集合的元素个数与所述已完成工具调用的总次数；所述元数据事件不包含所述工具输出本身，从而诊断记录的体积不随工具输出的体积增长。}

\p{在一个实施方式中，当所述无新证据连续计数在步骤 S207 之后达到第一阈值时，所述循环控制单元 105 向本次的工具结果尾部追加策略变更指示。所述策略变更指示的文本指明：其由所述智能体运行时 100 依据工具结果生成而非由用户给出，最近若干次已完成的工具调用未增加可区分的证据，下一次动作必须改变假设、来源、目标或机制并说明其能够产生何种新的观察，若不存在此种动作则应停止并报告已获得的证据、重复出现的结果与尚未解决的缺口。所述策略变更指示被追加到工具结果之中而非作为一条新的用户消息插入，从而不改变消息角色的交替次序。在一个实施方式中，注入所述策略变更指示的同时记录当时的所述无新证据连续计数，并在所述无新证据连续计数相对于所记录的值的增量小于所述第一阈值之前不再次注入，以避免在同一段连续计数内反复注入相同的指示。}

\p{在一个实施方式中，所述第一阈值取 4。}

\cnexample{实施例三：执行前的阻断判定}

\p{参照图 3，当所述模型接口单元 101 接收到所述大语言模型 110 生成的一次新的工具调用后，在将其交由所述工具执行单元 102 执行之前，所述循环控制单元 105 执行如下步骤。}

\p{步骤 S301，判断当前是否处于所述第三模式。若否，则不作阻断判定，直接转入执行。}

\p{步骤 S302，判断该次工具调用的工具名称是否属于预设轮询类工具集合，并据此确定本次判定所用的连续计数阈值。若属于，则该阈值取所述第二阈值与预设倍数的乘积；否则取所述第二阈值。在一个实施方式中，所述预设轮询类工具集合以工具名称的匹配规则定义，所述匹配规则包括：工具名称等于表示等待的名称，或工具名称包含表示等待某一条件的子串、表示状态查询的子串、表示读取会话线程的子串或表示查询交接状态的子串。在一个实施方式中，所述第二阈值取 6，所述预设倍数取 2。}

\p{步骤 S303，判断所述无新证据连续计数是否达到步骤 S302 确定的阈值。若未达到，则不阻断，转入执行。}

\p{步骤 S304，按实施例二步骤 S201 所述规则生成该次工具调用的调用指纹，并从所述证据账本单元 104 取出该调用指纹对应的调用指纹计数与同结果计数；若不存在对应记录，则取零。}

\p{步骤 S305，判断所述调用指纹计数是否不小于第三阈值，且所述同结果计数是否不小于第四阈值。若任一条件不满足，则不阻断，转入执行。在一个实施方式中，所述第三阈值与所述第四阈值均取 3。}

\p{由步骤 S303 与步骤 S305 可见，阻断要求三个条件同时成立：本用户回合已连续若干次未获得新证据、该次工具调用此前已以完全相同的工具名称与入参发起过足够多次、且该工具调用此前已足够多次产生了同一个结果。因此，一次此前未曾发起过的工具调用、或一次虽曾发起但每次结果均不相同的工具调用，在所述无新证据连续计数已经很高时依然可以执行；改变来源、目标、检索词或机制的动作始终保持可用。}

\p{步骤 S306，三个条件同时成立时，所述循环控制单元 105 不再将该次工具调用交由所述工具执行单元 102，而是合成一条拒绝结果作为其工具结果写入所述会话上下文。所述拒绝结果为结构化对象，含有取值表示已阻断的状态字段、取值表示无新证据循环的原因字段、说明文本字段以及分别记录所述无新证据连续计数、所述调用指纹计数与所述同结果计数的三个数值字段；所述说明文本指明该工具此前已以相同入参运行的次数、产生相同结果的次数与当前的连续计数，并要求在尝试下一个工具之前改变假设、来源、目标或机制，或者停止并报告尚未解决的缺口。所述拒绝结果被标记为失败，从而在所述会话上下文、用户界面与诊断记录中不被呈现为一次成功的工具调用。}

\p{步骤 S307，所述拒绝结果不参与实施例二所述的账本更新。具体地，在执行账本更新之前先判断工具输出中是否含有表示无新证据循环的原因字段，若含有则直接返回而不更新任何计数与集合。其原因在于，所述拒绝结果描述的是账本自身的状态，并非从被请求的工具处取得的观察；若将其计入账本，其作为一个此前未出现过的结果会把所述无新证据连续计数清零，从而使下一次相同的工具调用重新获得执行机会，阻断条件将永远无法稳定成立。}

\p{在一个实施方式中，当所述大语言模型 110 在一次响应中生成多个并行的工具调用时，所述循环控制单元 105 按其在响应中的次序逐个执行上述判定，账本的更新仅依据已经完成的工具结果进行，并且已经派发执行的工具调用不因同一响应中后续的判定结果而被撤销。}

\cnexample{实施例四：大体积工具输出的外置与证据身份桥接}

\p{参照图 4，本实施例说明超过外置阈值的工具输出如何被移出会话上下文，以及如何使该外置不破坏实施例二与实施例三所述的判定。}

\p{步骤 S401，所述智能体运行时 100 确定本次的外置阈值。具体地，取本用户回合当前配置的上下文压缩令牌预算，按预设的每令牌字符数将其折算为字符数，取该字符数的预设比例，并将所得结果钳制于预设下界与预设上界之间。在一个实施方式中，所述每令牌字符数取 4，所述预设比例取百分之五，所述预设下界取 8000 个字符，所述预设上界取 32000 个字符。由此，上下文窗口较小的模型获得较小的外置阈值，上下文窗口较大的模型获得较大的外置阈值，而所述钳制使两端的取值保持可预期。在另一个实施方式中，运行于桌面应用中的所述智能体运行时 100 采用固定的外置阈值，例如 64000 个字符。}

\p{步骤 S402，判断该次工具调用的工具输出原文的字符数与其在会话上下文中的呈现文本的字符数是否均不超过所述外置阈值。若均不超过，则不作外置，仅对呈现文本按既有的结果长度上限作截断，流程结束。}

\p{步骤 S403，所述产物存储单元 106 对所述工具输出原文计算内容摘要值。}

\p{步骤 S404，判断是否可以复用已存在的产物文件。具体地，尝试将本次的呈现文本解析为结构化对象，并在其中递归查找存储路径字段；若查找到取值非空的存储路径，则对该路径求规范化的绝对路径，并判断其是否为一个已存在的文件且位于预设产物目录之内。仅在上述条件全部满足时复用该产物文件，不再次写入。所述目录归属判断使得会话上下文中出现的任意路径不能诱导所述产物存储单元 106 将产物目录之外的文件当作产物复用。}

\p{步骤 S405，若不能复用，则生成产物文件的文件名并写入。所述文件名由所述内容摘要值的前若干位、经字符过滤后的所述工具名称与经字符过滤后的所述调用标识依次以分隔符拼接并附加扩展名构成；所述字符过滤将非字母、非数字且非连字符、非下划线的字符替换为下划线并截断至预设长度，过滤后为空时以预设的固定字符串代替。随后，仅在该文件尚不存在时，以原子替换的方式将所述工具输出原文写入该文件。由此，同一份内容在其文件已经存在时不会被重复写入，而不同调用产生的同名冲突被所述调用标识消除。}

\p{步骤 S406，所述产物存储单元 106 以引用式结果替换所述会话上下文中的该工具输出。所述引用式结果为结构化对象，含有：取值表示已引用的状态字段、预览字段、存储路径字段、存储体积字段、原始字符数字段、内容摘要值字段与检索提示字段。所述预览由所述呈现文本按预设规则截取得到：取预设预览上限与所述外置阈值中的较小者为预览长度上限，在所述呈现文本的字符数不超过该上限时直接取其全文，否则从该上限中扣除省略标记本身的长度后，将剩余额度按预设比例分配给首部与尾部，并由首部片段、所述省略标记与尾部片段依次拼接而成；所述省略标记声明被省略的是该引用式结果的中部及其原始字符数。在一个实施方式中，所述预设预览上限取 12000 个字符，所述首部与尾部的分配比例取四分之三与四分之一。所述检索提示指明：应对所述存储路径以较窄的行范围读取，或以具体的模式对其检索，非必要不整体重新载入。}

\p{步骤 S407，证据身份的桥接。所述产物文件的存储路径中含有所述调用标识，因而同一份内容在不同次工具调用下必然落在不同的路径上。为此，实施例二步骤 S204 在识别到引用式结果时，以其中的内容摘要值字段的取值而非存储路径字段的取值作为规范化结果的生成依据。由此，两次工具调用只要取回了完全相同的内容，其在会话上下文中的引用式结果虽然路径不同，所生成的证据指纹仍然相同，判定结果为“重复”而非“新”。}

\p{在一个实施方式中，所述智能体运行时 100 在生成证据指纹时使用工具输出原文，而在会话上下文中写入引用式结果；当所述会话上下文其后被压缩并按实施例五重建账本时，可用的只有已写入的引用式结果，此时步骤 S204 所述的桥接使得重建得到的判定与外置之前保持一致的重复关系。若不作此桥接，则所述重建会把同一份内容的每一个副本都当作新证据，仅仅因为它被写在了另一个文件名之下，所述无新证据连续计数将在重建后被清零，实施例三所述阻断在压缩之后即告失效。}

\p{在一个实施方式中，当所述工具调用返回的是图像类结果时，不执行本实施例所述外置，以免破坏该结果既有的呈现形式。}

\cnexample{实施例五：上下文压缩时的账本重建}

\p{参照图 5，当所述上下文压缩单元 108 判定所述会话上下文超出预算而需要压缩时，执行如下步骤。}

\p{步骤 S501，确定重建的选取范围。自所述会话上下文的末尾向前查找最后一条角色为用户且不是由所述智能体运行时 100 内部生成的用户消息，以该消息之后的位置为起点，以所述会话上下文的末尾为终点。由此，重建仅覆盖当前用户回合，而运行时内部生成的用户消息（例如压缩延续消息或内部提示消息）不会被误认为新回合的起点从而截断范围。}

\p{步骤 S502，遍历所述选取范围内的每条消息的每个内容块，将遇到的工具调用块以其调用标识为键记录其工具名称与调用入参。}

\p{步骤 S503，对遇到的每个工具结果块，以其调用标识在步骤 S502 所记录的内容中取回对应的调用入参；若未取回，则以预设的空结构化对象字符串代替。随后以该工具结果块的工具名称、取回的调用入参、该工具结果块的工具输出与失败标记为参数，读取观察记录中对应的规范化规则版本，按实施例二步骤 S201 至步骤 S208 所述的完全相同的规则更新一个新建的账本。对于引用式结果，优先读取其中的内容摘要值作为证据身份；对于合成拒绝结果，依据其合成结果标记跳过账本更新。}

\p{步骤 S504，遍历结束后，所述新建的账本即为重建结果，其所含的调用指纹计数、证据指纹集合、同结果计数与无新证据连续计数与压缩之前保持一致的判定关系。由于全部规则均为确定性的规范化、摘要和结果类型判定运算，重建不需要调用所述大语言模型 110 或外部工具，也不引入额外的令牌或网络请求消耗。}

\p{步骤 S505，所述上下文压缩单元 108 由所述重建结果导出聚合事实并写入压缩产物。所述聚合事实包括：所述证据指纹集合的元素个数、已完成工具调用的总次数、其中重复的次数与空结果的次数；当所述无新证据连续计数大于零时，还包括当前的连续计数长度。所述聚合事实被写入压缩后的摘要中的专门小节，并同时写入跨压缩保留的固定上下文，从而恢复后的会话即使不再持有原始消息，也仍然知悉此前的工作处于重复而非推进的状态。}

\p{在一个实施方式中，所述上下文压缩单元 108 同时收集所述会话上下文中出现过的引用式结果的存储路径并写入压缩产物，使得压缩之后所述大语言模型 110 仍可按需读取被外置的工具输出原文。}

\cnexample{实施例六：变形与等同替换}

\p{在一个实施方式中，所述摘要值由密码学散列函数计算得到并以十六进制字符串表示。本领域技术人员可以理解，任何具有足够抗碰撞性的确定性映射均可用于生成所述调用指纹与所述证据指纹，所述内容摘要值亦然。}

\p{在一个实施方式中，所述第一阈值、所述第二阈值、所述第三阈值、所述第四阈值与所述预设倍数为可配置参数。所述第一阈值小于所述第二阈值；所述第三阈值与所述第四阈值不小于二。}

\p{在一个实施方式中，所述预设轮询类工具集合以显式列表给出，或由工具的声明属性给出，而不限于按工具名称匹配。}

\p{在一个实施方式中，所述产物文件存储于项目本地的临时目录之下；在另一个实施方式中，所述产物文件存储于所述智能体运行时 100 可访问的对象存储之中，此时所述存储路径为对应的对象地址，所述桥接规则不变。}

\p{在一个实施方式中，所述方法以计算机程序的形式存储于计算机可读存储介质中，由处理器读取并执行；所述计算机可读存储介质包括只读存储器、随机存取存储器、磁盘或光盘。}

\p{在一个实施方式中，所述智能体运行时 100 的各单元可以合并或拆分，例如所述指纹生成单元 103 可以并入所述证据账本单元 104，所述循环控制单元 105 可以并入所述模型接口单元 101，只要其共同实现前述各步骤即可。}

\p{以上所述仅为本发明的较佳实施例，并不用以限制本发明。凡在本发明的精神和原则之内所作的任何修改、等同替换和改进等，均应包含在本发明的保护范围之内。}

\end{document}
